\makeatletter
\@ifundefined{examheaderloaded}{%
\gdef\examheaderloaded{}%
\usepackage{xcolor}
\usepackage{tikz}
\usepackage{tabularx}
\usepackage{array}
\usepackage{enumitem}
\usepackage{multicol}
\usepackage{fancyvrb}
\usepackage{needspace}
\usepackage[most]{tcolorbox}
\usepackage{fancyhdr}
\usepackage{etoolbox}
\usepackage{xstring}
\usepackage{pdflscape}
\usetikzlibrary{calc}

\definecolor{examprimary}{HTML}{1565C0}
\definecolor{examprimarylight}{HTML}{DCE8F8}
\definecolor{examink}{HTML}{172033}
\definecolor{examlightgray}{HTML}{F5F5F5}
\definecolor{exammidgray}{HTML}{CCCCCC}
\definecolor{examrulegray}{HTML}{AAAAAA}
\definecolor{examwarn}{HTML}{B45309}
\definecolor{examtip}{HTML}{2E7D32}
\definecolor{examcardback}{HTML}{F4F8FD}
\newcommand{\headingfont}{\fontfamily{qpl}\selectfont}
\newcommand{\uhr}{%
  \tikz[baseline=-0.4ex, line width=0.7]{\draw (0,0) circle (0.42em); \draw (0,0) -- (0,0.3em); \draw (0,0) -- (0.24em,0);}}
\ifPDFTeX\DeclareUnicodeCharacter{25A1}{\ensuremath{\square}}\fi
\ifPDFTeX\DeclareUnicodeCharacter{2610}{\ensuremath{\square}}\fi
\ifPDFTeX\DeclareUnicodeCharacter{2713}{\ensuremath{\checkmark}}\fi
\ifPDFTeX\DeclareUnicodeCharacter{25B6}{\ensuremath{\blacktriangleright}}\fi
\ifPDFTeX\DeclareUnicodeCharacter{2022}{\ensuremath{\bullet}}\fi

\AtBeginDocument{%
  \color{examink}%
  \setlength{\parindent}{0pt}%
  \setlength{\parskip}{0.45em}%
  \renewcommand{\arraystretch}{1.18}%
}

\pagestyle{fancy}
\fancyhf{}
\makeatletter
\newif\ifexamshownameheader
\examshownameheadertrue
\newcommand{\loesungsdokument}{\examshownameheaderfalse}
\newcommand{\marksolutiondocument}[1]{%
  \IfSubStr{#1}{Musterl}{\loesungsdokument}{}%
}
\newcommand{\examfollowpagegeometry}{%
  \newgeometry{a4paper,top=14mm,right=13mm,bottom=5mm,left=14mm,headsep=3mm}%
}
\newcommand{\examdefaultpagegeometry}{%
  \newgeometry{a4paper,top=14mm,right=13mm,bottom=5mm,left=14mm,headsep=3mm}%
}
\newcommand{\exampagebreak}{%
  \newpage
  \ifodd\value{page}%
    \examdefaultpagegeometry
  \else
    \examfollowpagegeometry
  \fi
}
\fancyhead[C]{\ifdefined\studentfieldsinheader\ifnum\value{page}=1\studentheaderrow\fi\fi}
\fancyhead[L]{\ifexamshownameheader\ifnum\value{page}>1\ifodd\value{page}\else\footnotesize\textbf{Name:}\ \rule{8.5cm}{0.4pt}\fi\fi\fi}
\fancyhead[R]{\ifnum\value{page}>1\ifodd\value{page}\else\footnotesize\color{gray}\@title\ $\cdot$ Seite \thepage\fi\fi}
\makeatother
\renewcommand{\headrulewidth}{0pt}
\setlength{\headheight}{13pt}

\newcommand{\coverheader}[3]{%
  \marksolutiondocument{#1}%
  \begin{center}
    {\headingfont\Large\bfseries\color{examprimary}#1\par}
    \vspace{1mm}
    {\large #2\par}
    {\small\color{gray}#3\par}
  \end{center}
  \vspace{0.45cm}
}

\newcommand{\fieldline}[1][4cm]{\rule{#1}{0.4pt}}
\newcommand{\studentheaderrow}{%
  \begingroup\normalfont\fontsize{9.5}{11}\selectfont\color{black}%
  \begin{tabularx}{\headwidth}{@{}lXl p{40mm}@{}}
  \textbf{Name:} & \fieldline[\linewidth] & \textbf{Datum:} & \fieldline[\linewidth]
  \end{tabularx}%
  \endgroup
}
\newcommand{\studentfields}{%
  \ifdefined\studentfieldsinheader\else
  \noindent\begin{tabularx}{\linewidth}{@{}lXlX@{}}
  \textbf{Name:} & \fieldline[\linewidth] & \textbf{Datum:} & \fieldline[\linewidth]
  \end{tabularx}
  \vspace{0.6cm}
  \fi
}
\newcommand{\gradingfields}[1]{%
  \noindent\begin{tabularx}{\linewidth}{@{}lXllX@{}}
  \textbf{Punkte:} & \fieldline[\linewidth] & / #1 & \textbf{Notenpunkte:} & \fieldline[\linewidth]\\[0.45cm]
  \textbf{Zeit:} & 90 Min & & \textbf{Hilfsmittel:} & keine
  \end{tabularx}
  \vspace{0.6cm}
}

\newsavebox{\miniabbox}
\newenvironment{miniab}{%
  \par
  \begin{lrbox}{\miniabbox}%
  \begin{minipage}[t]{\linewidth}%
  \definecolor{minianswergray}{HTML}{EEEEEE}%
  \renewcommand{\antwortkasten}[1]{%
    \par\vspace{0.2em}\noindent
    \begin{tikzpicture}
      \fill[minianswergray,rounded corners=2pt] (0,0) rectangle (\linewidth,##1);
    \end{tikzpicture}\par\vspace{0.35em}%
  }%
  \setlength{\parskip}{0.2em}%
  \renewcommand{\arraystretch}{1.05}%
  \setlength{\LTpre}{2pt}%
  \setlength{\LTpost}{2pt}%
  \renewcommand{\toprule}{}%
  \renewcommand{\midrule}{}%
  \renewcommand{\bottomrule}{}%
  \renewcommand{\antwortlinie}[1][4.5cm]{\textcolor{examrulegray}{\rule{##1}{0.4pt}}}%
  \renewcommand{\studentfields}{%
    \noindent{\bfseries Name:}\ \textcolor{examrulegray}{\fieldline[6cm]}\hfill{\bfseries Datum:}\ \textcolor{examrulegray}{\fieldline[2.5cm]}\par
    \vspace{0.15cm}%
  }%
  \renewenvironment{inputtable}[1][0pt]{%
    \begingroup\renewcommand{\arraystretch}{1.15}\setlength{\extrarowheight}{\dimexpr##1+1pt\relax}%
  }{\endgroup}%
  \renewcommand\section{\@startsection{section}{1}{\z@}{1.1em}{0.65em}{\normalfont\bfseries\color{examprimary}}}%
  \renewcommand\subsection{\@startsection{subsection}{2}{\z@}{1.1em}{0.65em}{\normalfont\bfseries\color{examprimary}}}%
}{%
  \end{minipage}%
  \end{lrbox}%
  \noindent\usebox{\miniabbox}\par
  \vfill
  \noindent{\color{exammidgray}\dotfill}\par
  \vfill
  \noindent\usebox{\miniabbox}\par
}

\newenvironment{hinweis}{\begingroup\small}{\par\endgroup}

\newtcolorbox{infobox}[1]{%
  enhanced, breakable,
  colback=examlightgray,
  colframe=examprimary,
  colbacktitle=examprimary,
  coltitle=white,
  fonttitle=\headingfont\bfseries\small,
  title={#1},
  boxrule=0.5pt, arc=2pt,
  left=6pt, right=6pt, top=3pt, bottom=3pt,
  before skip=0.5em, after skip=0.6em
}

\newtcolorbox{warnbox}[1]{%
  enhanced, breakable,
  colback=examlightgray,
  colframe=examwarn,
  colbacktitle=examwarn,
  coltitle=white,
  fonttitle=\headingfont\bfseries\small,
  title={#1},
  boxrule=0.5pt, arc=2pt,
  left=6pt, right=6pt, top=3pt, bottom=3pt,
  before skip=0.5em, after skip=0.6em
}

\newtcolorbox{tippbox}[1]{%
  enhanced, breakable,
  colback=examlightgray,
  colframe=examtip,
  colbacktitle=examtip,
  coltitle=white,
  fonttitle=\headingfont\bfseries\small,
  title={#1},
  boxrule=0.5pt, arc=2pt,
  left=6pt, right=6pt, top=3pt, bottom=3pt,
  before skip=0.5em, after skip=0.6em
}

\newenvironment{aufgabe}[3]{%
  \par\vspace{0.45em}
  \Needspace{8cm}
  \noindent{\headingfont\large\bfseries\color{examprimary}Aufgabe #1 - #2\hfill{\normalsize \ifstrempty{#3}{}{#3 Punkte}}}\par
  \vspace{0.05em}
}{\par\vspace{0.25em}}

\newcommand{\fortsetzung}[2]{%
  \par\vspace{0.2em}
  \noindent{\headingfont\bfseries\color{examprimary}Aufgabe #1 - #2 (Fortsetzung)}\par
}

\newcommand{\currentteilpunkte}{}
\newenvironment{teil}[2]{%
  \renewcommand{\currentteilpunkte}{#2}%
  \par\vspace{0.25em}\noindent
  \textbf{#1)}\hspace{0.45em}%
  \begin{minipage}[t]{0.82\linewidth}\ignorespaces
}{%
  \end{minipage}\hfill{\small\color{gray}(\currentteilpunkte\ P)}\par
}

\newcommand{\antwortkasten}[1]{%
  \par\vspace{0.2em}\noindent
  \begin{tikzpicture}
    \draw[fill=examlightgray,draw=exammidgray,line width=0.4pt,rounded corners=2pt] (0,0) rectangle (\linewidth,#1);
  \end{tikzpicture}
  \par\vspace{0.35em}
}

\newcommand{\answergrid}[3]{%
  \par\vspace{0.35em}\noindent
  \begin{tikzpicture}
    \fill[examlightgray,rounded corners=2pt] (0,0) rectangle ({#3*#1},{#2*#1});
    \draw[rounded corners=2pt,draw=exammidgray,line width=0.4pt] (0,0) rectangle ({#3*#1},{#2*#1});
    \draw[step=#1,draw=examrulegray,line width=0.25pt] (0,0) grid ({#3*#1},{#2*#1});
  \end{tikzpicture}
  \par\vspace{0.45em}
}

\newcommand{\antwortlinie}[1][4.5cm]{\rule{#1}{0.4pt}}

\definecolor{examanswerfield}{HTML}{EEEEEE}
\newcommand{\platz}[1][1.5cm]{\colorbox{examanswerfield}{\rule[-0.9mm]{0pt}{0.45cm}\makebox[#1][c]{}}}

\DefineVerbatimEnvironment{Highlighting}{Verbatim}{commandchars=\\\{\},fontsize=\small}
\newenvironment{Shaded}{\par\smallskip}{\smallskip}
\newcommand{\NormalTok}[1]{#1}
\newcommand{\KeywordTok}[1]{\textbf{\textcolor{examprimary}{#1}}}
\newcommand{\ControlFlowTok}[1]{\textbf{\textcolor{examprimary}{#1}}}
\newcommand{\ImportTok}[1]{\textbf{\textcolor{examprimary}{#1}}}
\newcommand{\BuiltInTok}[1]{\textcolor{purple}{#1}}
\newcommand{\OperatorTok}[1]{#1}
\newcommand{\DecValTok}[1]{\textcolor{brown}{#1}}
\newcommand{\BaseNTok}[1]{\textcolor{brown}{#1}}
\newcommand{\FloatTok}[1]{\textcolor{brown}{#1}}
\newcommand{\CharTok}[1]{\textcolor{teal}{#1}}
\newcommand{\StringTok}[1]{\textcolor{teal}{#1}}
\newcommand{\CommentTok}[1]{\textit{\textcolor{gray}{#1}}}
\newcommand{\OtherTok}[1]{#1}
\newcommand{\FunctionTok}[1]{\textcolor{purple}{#1}}
\newcommand{\VariableTok}[1]{#1}
\newcommand{\AttributeTok}[1]{#1}
\newcommand{\DataTypeTok}[1]{\textcolor{purple}{#1}}
\newcommand{\ConstantTok}[1]{\textcolor{brown}{#1}}
\newcommand{\SpecialCharTok}[1]{\textcolor{teal}{#1}}
\newcommand{\VerbatimStringTok}[1]{\textcolor{teal}{#1}}
\newcommand{\SpecialStringTok}[1]{\textcolor{teal}{#1}}
\newcommand{\ErrorTok}[1]{\textcolor{red}{#1}}
\newcommand{\AlertTok}[1]{\textbf{\textcolor{red}{#1}}}
\newcommand{\InformationTok}[1]{\textcolor{gray}{#1}}
\newcommand{\WarningTok}[1]{\textbf{\textcolor{red}{#1}}}

\newenvironment{codehinweis}{\begingroup\small\itshape\color{gray}}{\par\endgroup}
\newenvironment{referenz}{\begingroup\footnotesize}{\par\endgroup}
\newenvironment{wahrfalschlayout}{\begingroup\renewcommand{\arraystretch}{1.55}}{\endgroup}
\newenvironment{inputtable}[1][0pt]{\begingroup\renewcommand{\arraystretch}{1.45}\setlength{\extrarowheight}{#1}}{\endgroup}
\newenvironment{prooftable}{\begingroup\renewcommand{\arraystretch}{1.6}}{\endgroup}
\newenvironment{columnstwo}{\begin{multicols}{2}}{\end{multicols}}

\newenvironment{kartensatz}[1]{%
  \par\vspace{0.4em}%
  \ifstrempty{#1}{}{{\headingfont\large\bfseries\color{examprimary}#1}\par\vspace{0.15em}}%
}{%
  \par\vspace{0.2em}%
}
\newcommand{\kartezelle}[2]{%
  \begin{tcolorbox}[enhanced, colback=examcardback, colframe=examprimary,
    boxrule=0.8pt, arc=4pt, left=4pt, right=4pt, top=3pt, bottom=3pt,
    before skip=0pt, after skip=0pt, height=17mm, valign=center, halign=center]
    \ifstrempty{#1}{}{%
      {\setlength{\fboxsep}{1.5pt}\colorbox{examprimary}{\textcolor{white}{\footnotesize\textbf{#1}}}}%
      \par\vspace{2pt}%
    }%
    {\normalsize #2}%
  \end{tcolorbox}%
}

\newcommand{\notenspiegeltable}{%
  \par\small
  \begin{tabularx}{\linewidth}{|X|*{17}{>{\centering\arraybackslash}p{0.037\linewidth}|}}
  \hline
  \textbf{NP} & \textbf{15} & \textbf{14} & \textbf{13} & \textbf{12} & \textbf{11} & \textbf{10} & \textbf{9} & \textbf{8} & \textbf{7} & \textbf{6} & \textbf{5} & \textbf{4} & \textbf{3} & \textbf{2} & \textbf{1} & \textbf{0} & \textbf{ø}\\
  \hline
  \textbf{Anzahl} &&&&&&&&&&&&&&&&&\\
  \hline
  \end{tabularx}
  \normalsize\par
}

\newenvironment{erwartungshorizont}[1]{%
  \clearpage
  \begin{landscape}
  \thispagestyle{empty}
  \par\vspace{0.4em}
  \noindent{\headingfont\large\bfseries\color{examprimary}#1}\par
  \vspace{0.1em}
  \begingroup\small
}{\par\endgroup\end{landscape}}

\newenvironment{erwartungstabelle}{%
  \par\vspace{0.2em}
  \renewcommand{\arraystretch}{1.25}%
  \setlength{\tabcolsep}{4pt}%
  \begin{longtable}{|>{\raggedright\arraybackslash}p{0.49\linewidth}|>{\raggedright\arraybackslash}p{0.35\linewidth}|>{\centering\arraybackslash}p{2cm}|}
  \hline
  \textbf{Erwartete Lösung} & \textbf{Schülerlösung} & \textbf{Punkte / erreichte Punkte}\\
  \hline
  \endfirsthead
  \hline
  \textbf{Erwartete Lösung} & \textbf{Schülerlösung} & \textbf{Punkte / erreichte Punkte}\\
  \hline
  \endhead
}{%
  \end{longtable}
}

\long\def\ehzeile#1#2{%
  #1 & \rule[-5cm]{0pt}{5cm} & #2 / \\
  \hline
}

\newenvironment{loesung}[1]{%
  \par\vspace{0.25em}
  {\bfseries\color{examprimary}Aufgabe #1}\par
  \begingroup\small
}{\par\endgroup}

\newenvironment{versionen}{\par\noindent}{\par}
\newenvironment{versionleft}{\begin{minipage}[t]{0.48\linewidth}}{\end{minipage}\hfill}
\newenvironment{versionright}{\begin{minipage}[t]{0.48\linewidth}}{\end{minipage}}

\newenvironment{mathedeckblatt}[2]{%
  \marksolutiondocument{#1}%
  \begin{center}
  {\headingfont\Large\bfseries\color{examprimary}#1\par}
  \vspace{-3pt}
  {\color{examprimary}\rule{0.55\linewidth}{0.55pt}\par}
  \ifexamshownameheader
    \vspace{0.8cm}
    \begin{tabularx}{\linewidth}{@{}Xr@{}}
    Name: \fieldline[6.4cm] & Datum: #2\fieldline[3cm]
    \end{tabularx}
  \fi
  \end{center}
  \vspace{0.3cm}
}{}

\newcommand{\notetable}{%
  \par
  \begin{tabularx}{\linewidth}{|X|*{7}{>{\centering\arraybackslash}X|}}
  \hline
  Note & 1 & 2 & 3 & 4 & 5 & 6 & ø\\
  \hline
  Anzahl &&&&&&&\\
  \hline
  Notw. Punkte &&&&&&&\\
  \hline
  \end{tabularx}
  \par
}

\newcommand{\foerderplankopf}[6]{%
  {\Large\bfseries Christian-Wirth-Schule}\hfill
  \includegraphics[width=3.4cm]{assets/cws-logo.png}\par
  {\small Schloßplatz 1 · 61250 Usingen}\par
  \rule{\linewidth}{0.6pt}\par
  \vspace{0.2cm}
  \begin{tabularx}{\linewidth}{|X|X|X|}
  \hline
  \textbf{Förderplan für:} #1 & \textbf{Klasse:} #2 & \textbf{Schuljahr:} #3\\
  \hline
  \textbf{Verfasser:} #4 & \textbf{Fach:} #5 & \textbf{Datum:} #6\\
  \hline
  \end{tabularx}
}

\newenvironment{foerderbereich}[1]{%
  \par\vspace{0.45cm}
  \noindent{\headingfont\large\bfseries Förderbereich: #1}\par
  \rule{\linewidth}{0.5pt}\par
  \begingroup\renewcommand{\arraystretch}{1.3}
}{\endgroup}

\newcommand{\unterschriften}{%
  \par\vspace{1.2cm}
  \begin{tabularx}{\linewidth}{XXX}
  \rule{\linewidth}{0.4pt} & \rule{\linewidth}{0.4pt} & \rule{\linewidth}{0.4pt}\\
  {\footnotesize Ort, Datum, Unterschrift}\newline\textbf{Lehrer/in} &
  {\footnotesize Ort, Datum, Unterschrift}\newline\textbf{Erziehungsberechtigte/r} &
  {\footnotesize Ort, Datum, Unterschrift}\newline\textbf{Schüler/in}
  \end{tabularx}
}
}{}
\makeatother
